% Source: https://artofproblemsolving.com/community/c3h2049608_a_hard_graph
% Stable identity: aops-topic-2049608-q01
The equation holds exactly when $x+y-2\ge0$ and the squares are equal, since the square root is nonnegative, and the condition $x^2+y^2-4\ge0$ then follows automatically from the squared equation. Squaring:
\[x^2+y^2-4=x^2+y^2+4+2xy-4x-4y\iff 2xy-4x-4y+8=0\iff(x-2)(y-2)=0 .\]
So $x=2$ or $y=2$, together with $x+y\ge2$:
\begin{itemize}
\item $x=2$ requires $y\ge0$;
\item $y=2$ requires $x\ge0$.
\end{itemize}

\begin{asy}
usepackage("amsmath");
size(6cm);
defaultpen(fontsize(10pt));
draw((-1,0)--(6,0),gray,Arrow);draw((0,-1)--(0,6),gray,Arrow);
label("$x$",(6,0),E);label("$y$",(0,6),N);
draw((2,0)--(2,5.6),blue+linewidth(1.2));
draw((0,2)--(5.6,2),blue+linewidth(1.2));
draw((2,-1)--(2,0),blue+dashed);draw((-1,2)--(0,2),blue+dashed);
draw((-1,3)--(3,-1),red+dashed);
label("$x+y=2$",(3,-1),E,red);
dot((2,0),blue);dot((0,2),blue);
label("$(2,0)$",(2,0),SE);label("$(0,2)$",(0,2),NE);
label("$x=2$",(2,5.6),E,blue);label("$y=2$",(5.6,2),N,blue);
label("\mbox{}",(min(currentpicture,user=true)+max(currentpicture,user=true))/2,black); // keeps the text after the figure black when LaTeX typesets the labels inline
\end{asy}

The graph is the union of two rays:
\[\boxed{\{(2,y):y\ge0\}\ \cup\ \{(x,2):x\ge0\}},\]
the ray going up from $(2,0)$ along $x=2$ and the ray going right from $(0,2)$ along $y=2$, which cross at $(2,2)$. The dashed parts of the lines $x=2$ and $y=2$ are extraneous: there $x+y-2<0$.
